\uuid{3REh}
\titre{Noyau/Image/Injection/Surjection/Théorème du rang}
\auteur{hanani ; mzari}
\organisation{exo7}
\licence{CC BY-NC-SA 4.0}
\datecreate{2020-01-20}
\niveau{L1}
\module{Algèbre}
\chapitre{Application linéaire}
\sousChapitre{Définition}
\difficulte{3}
\theme{Noyau et image, Noyau, Image, Injection, Surjection, Théorème du rang}
\exo7qcmid{lille-2.applications-lineaires.025}

\contenu{
\begin{question}
\qclassification{107.01}
\qtags{motcle=Noyau/Image/Injection/Surjection/Théorème du rang}
Soit $E$ et $F$ deux $\Rr$-espaces vectoriels de dimensions finies
et $f$ une application linéaire de $E$ dans $F$. Quelles sont les assertions vraies ?
\begin{answers}
\bad{Si $f$ est injective, alors $f$ est surjective}.
\bad{Si $f$ est surjective, alors $f$ est injective}.
\bad{Si $\dim E = \dim F$, alors $f$ est bijective}.
\good{Si $f$ est bijective, alors $\dim E = \dim F$}.
\end{answers}
\begin{explanations} Si $\dim E \neq \dim F$, alors $f$ peut-être injective (resp. surjective) sans qu'elle soit surjective (resp. injective).
\vskip0mm
Si $\dim E = \dim F$, $f$ n'est pas nécessairement bijective.
\vskip0mm
Par contre, si  $f$ est bijective, comme  $E$ et $F$ sont deux $\Rr$-espaces vectoriels de  dimensions finies, du théorème du rang, on déduit que $\dim E = \dim F$.
\end{explanations}
\end{question}
}
